%! AP Statistics — Unit 2, Lesson 2.9 — Warm-Up
\documentclass[10pt]{article}
\usepackage{apstats-article}
\usepackage{apstats-boxes}

\begin{document}

\pageheader{Unit 2, Lesson 2.9}{Warm-Up}
\namedateperiod

\vspace{0.3in}

% ── PART 1 ────────────────────────────────────────────────────────────────────
{\large\bfseries\color{navy} Part 1 \quad Residual Plots and Model Fit (Lessons 2.7--2.8)}

\vspace{0.12in}
\noindent Recall: a residual plot that shows a \textbf{curved pattern} is evidence that a
linear model is \textbf{not} appropriate for the data.

\vspace{0.1in}

\begin{enumerate}[leftmargin=1.6em, itemsep=0.22in]

  \item A regression of $y$ on $x$ is performed. The residual plot shows a clear U-shaped
        (parabolic) pattern.

        \begin{enumerate}[label=\textbf{(\alph*)}, leftmargin=1.6em, itemsep=8pt]
          \item What does the curved residual plot tell you about the linear model?\\[8pt]
                \writeline\\[1pt]\writeline

          \item The value $r^2 = 0.81$ for this fit. Does a high $r^2$ guarantee the linear
                model is appropriate? Explain.\\[8pt]
                \writeline\\[1pt]\writeline
        \end{enumerate}

  \item Given $r = 0.92$, $s_y = 4.8$, $s_x = 1.2$, $\bar{x} = 5$, $\bar{y} = 22$.

        \vspace{4pt}
        Compute the slope: $b = r\dfrac{s_y}{s_x} = $ \blank{4cm}

        \vspace{6pt}
        Compute the intercept: $a = \bar{y} - b\bar{x} = $ \blank{4cm}

\end{enumerate}

\vspace{0.28in}

% ── PART 2 ────────────────────────────────────────────────────────────────────
{\large\bfseries\color{navy} Part 2 \quad Logarithm Properties (Math Review)}

\vspace{0.12in}

\begin{tcolorbox}[colback=greenbg, colframe=greenacc, boxrule=0.8pt, arc=2mm,
    left=4mm, right=4mm, top=3mm, bottom=3mm]
    \small
    \textbf{Quick reminder} (base-10 log):\quad
    $\log(10^k) = k$ \quad $10^{\log x} = x$ \quad
    $\log(a \cdot b^x) = \log a + x\log b$
\end{tcolorbox}

\vspace{0.14in}

\begin{enumerate}[leftmargin=1.6em, itemsep=0.18in, resume]

  \item Evaluate without a calculator:

        \begin{enumerate}[label=\textbf{(\alph*)}, leftmargin=1.6em, itemsep=6pt]
          \item $\log(1000) = $ \blank{2.5cm}
          \item $10^{2.0} = $ \blank{2.5cm}
          \item If $\widehat{\log y} = 1.7 + 0.3x$ and $x = 2$,
                what is $\hat{y}$?\quad Show work: \blank{4cm}
        \end{enumerate}

  \item A scatterplot of $y$ vs.\ $x$ curves upward steeply. After applying $\log_{10}$ to
        the $y$-values the new scatterplot looks linear with $r = 0.97$.
        What does this suggest about the original relationship between $x$ and $y$?\\[8pt]
        \writeline\\[1pt]\writeline\\[1pt]\writeline

\end{enumerate}

\end{document}
